Systematic uncertainties arise from imperfect modelling of the detector response and from theoretical and modelling uncertainties affecting the simulated predictions of both background and signal processes.
They impact the analysis through variations in the overall normalisation and/or the shape of the final discriminant distributions.
All uncertainties are implemented as nuisance parameters in the profile-likelihood fit and are propagated coherently across analysis regions where appropriate.

% =========================================================
\subsection{Experimental uncertainties}
\label{sec:syst_exp}

\subsubsection{Integrated luminosity}

The uncertainty on the integrated luminosity of the full Run~2 dataset (2015--2018) amounts to 1.7\%~\cite{ATLAS-CONF-2019-021}.
It is derived from the primary luminosity calibration using the LUCID-2 detector~\cite{Avoni_2018} and is applied as a normalisation uncertainty to all simulated samples.

\subsubsection{Pile-up modelling}

The modelling of additional $pp$ interactions (pile-up) in simulation is corrected by reweighting the distribution of the average number of interactions per bunch crossing to match that observed in data.
The associated uncertainty is evaluated by varying the pile-up reweighting procedure within its recommended uncertainties and propagating the resulting variations to the final distributions.

\subsubsection{Lepton-related uncertainties}

Residual differences between data and simulation in lepton reconstruction, identification, isolation and trigger efficiencies are corrected using scale factors derived from tag-and-probe techniques~\cite{Aad:2019tso,Aad:2020gmm}.
Systematic uncertainties associated with these scale factors are propagated to the analysis discriminants.

For electrons, additional uncertainties related to the energy scale and resolution are considered, together with a dedicated uncertainty associated with charge mis-identification rejection.
For muons, uncertainties in the momentum scale and resolution are included, accounting for contributions from the inner detector and muon spectrometer.

\subsubsection{Small-$R$ jet uncertainties}

Jet-related uncertainties arise from JVT, the jet energy scale (JES) and the jet energy resolution (JER).

The JVT uncertainty is evaluated by varying the efficiency correction applied in simulation within its uncertainty.

The JES uncertainty is parameterised by multiple independent components that account for detector calibration effects, $\eta$ intercalibration, pile-up dependence, flavour composition and response, and additional effects relevant at high $p_{\mathrm{T}}$.
The JER uncertainty accounts for residual differences between data and simulation in the jet resolution modelling.
All jet-related variations are propagated to derived observables, including $E_{\mathrm{T}}^{\mathrm{miss}}$.

\subsubsection{$b$-tagging}
\label{b-tagging_uncertianties}

The $b$-tagging performance in simulation is corrected to match that in data using per-jet scale factors derived from dedicated calibration measurements.
The corresponding uncertainties are parameterised using eigenvector variations obtained from the diagonalisation of the calibration covariance matrices and are provided separately for $b$-jets and for the mis-tag rates of $c$- and light-flavour jets.
Additional uncertainties associated with the extrapolation of the calibration to the high-$p_{\mathrm{T}}$ regime are included following the recommended prescriptions~\cite{high_pt_extrapolation,btagging_rec_2023}.

\subsubsection{Missing transverse momentum}

The uncertainty on the missing transverse momentum, $E_{\mathrm{T}}^{\mathrm{miss}}$, is dominated by the modelling of the soft term reconstructed from tracks not associated with hard objects.
It is evaluated using data--MC comparisons of the transverse momentum balance between the hard-object and soft-term contributions and is parameterised by independent variations corresponding to the soft-term scale and resolution components.

% =========================================================
\subsection{Theory and modelling uncertainties}
\label{sec:syst_theory}

Theoretical and modelling uncertainties are assigned to the signal and simulated background processes.
These include variations of QCD renormalisation and factorisation scales, generator modelling differences (matrix element, parton shower and hadronisation), and additional normalisation uncertainties where theoretical precision is limited.

% ---------------------------------------------------------
\subsubsection{$t\bar{t}$ modelling uncertainties}

The modelling of $t\bar{t}$ production constitutes one of the dominant sources of systematic uncertainty in this analysis.

\paragraph{QCD scale variations.}

Uncertainties due to missing higher-order QCD corrections are evaluated by varying the renormalisation and factorisation scales in the nominal {\scshape Powheg}+ \allowbreak{\scshape Pythia8} configuration by factors of two and one-half.

\paragraph{Generator and parton shower modelling.}

The uncertainty associated with the choice of matrix-element generator is estimated by comparing the nominal {\scshape Powheg}+\allowbreak{\scshape Pythia8} configuration to an alternative {\scshape MadGraph5\allowbreak\_aMC@NLO}+\allowbreak{\scshape Pythia8} setup.
Parton shower and hadronisation uncertainties are evaluated by comparing {\scshape Powheg} predictions interfaced to {\scshape Pythia8} with those obtained using {\scshape Herwig7}.

\paragraph{Radiation modelling.}

Additional variations probe the modelling of initial- and final-state radiation and the choice of the $h_{\mathrm{damp}}$ parameter in \textsc{Powheg}.

\paragraph{Heavy-flavour normalisation.}

To account for residual mismodelling in the rates of $t\bar{t}$ production with additional heavy-flavour jets, a conservative normalisation uncertainty of 50\% is assigned independently to the $t\bar{t}+b$ and $t\bar{t}+c$ components.
A 5\% uncertainty is applied to the $t\bar{t}+\mathrm{light}$ component.

\paragraph{Four- vs five-flavour scheme comparison.}

An additional modelling uncertainty is derived from the comparison between the nominal 5FS \textsc{Powheg}+\textsc{Pythia8} prediction and an alternative 4FS NLO prediction obtained with \textsc{Powheg-Box}\allowbreak-\textsc{Res}.

\paragraph{NN reweighting.}

The uncertainty associated with the neural-network-based reweighting applied to the $t\bar{t}$ background is evaluated using dedicated correlated variations.

% ---------------------------------------------------------
\subsubsection{Other signal and background uncertainties}

\paragraph{\textbf{\boldmath Signal.}}

The signal modelling uncertainty includes renormalisation and factorisation scale variations evaluated in \textsc{MadGraph5\_aMC@NLO}, together with an additional PDF uncertainty.

\paragraph{\textbf{\boldmath $t\bar{t}t\bar{t}$.}}

Modelling uncertainties are derived from generator variations probing matrix-element scale choices and parton shower modelling.
An additional uncertainty on the inclusive $t\bar{t}t\bar{t}$ cross section is applied, amounting to $^{+7.8\%}_{-13.3\%}$~\cite{vanBeekveld:2022hty}.

\paragraph{\textbf{\boldmath $t\bar{t}W$, $t\bar{t}Z$, $t\bar{t}H$.}}

For $t\bar{t}W$ and $t\bar{t}Z$, generator modelling uncertainties are evaluated by comparing the nominal \textsc{MadGraph5\_aMC@NLO}+\textsc{Pythia8} prediction to alternative samples produced with \textsc{Sherpa}.
For $t\bar{t}H$, the nominal \textsc{Powheg}+\textsc{Pythia8} setup is compared to \textsc{MadGraph5\_aMC@NLO}+\textsc{Pythia8}.
Additional scale variations are included.
Normalisation uncertainties of 12\% and 10\% are applied to the inclusive $t\bar{t}Z$ and $t\bar{t}H$ cross sections, respectively.

\paragraph{\textbf{\boldmath Single top.}}

Modelling uncertainties are evaluated by comparing \textsc{Powheg} and \textsc{MadGraph5\allowbreak\_aMC@NLO} matrix-element generators, interfaced with \textsc{Pythia8} and \textsc{Herwig7}, and by varying QCD scales.
A conservative normalisation uncertainty of 30\% is assigned.

\paragraph{\textbf{\boldmath Minor processes.}}

A 60\% uncertainty is applied to $V+\mathrm{jets}$ production~\cite{Alwall:2007fs}.
For other minor backgrounds, a 50\% normalisation uncertainty is assigned.
